\clearpage
\VerdierPhysicalPageAnchor{0183}%
\begingroup\small\setlength{\parskip}{0pt}\setlength{\emergencystretch}{1.5em}
\setlength{\abovedisplayskip}{4pt plus 2pt minus 2pt}%
\setlength{\belowdisplayskip}{4pt plus 2pt minus 2pt}%

\VerdierSemanticAnchor{ch3d-1-3-2-morphism-m-cont}
态射 $m$ 是一个单态射，并且有一条正合列：

\VerdierSemanticAnchor{ch3d-1-3-2-exact-sequence}
\begin{equation}
0\longrightarrow c(\operatorname{id}_X)
\xrightarrow{m}c(\mu(S))[-1]
\longrightarrow c(\operatorname{id}_Z)[-1]
\longrightarrow0.
\tag{1.3.2.3}
\end{equation}
复形 $c(\operatorname{id}_X)$ 和 $c(\operatorname{id}_Z)$ 都同伦于零。因此，
复形 $c(\mu(S))$ 是无环的。故态射 $\mu(S)$ 是一个拟同构。

\VerdierSemanticAnchor{ch3d-1-3-2-k-diagram}
转到范畴 $K(\mathcal A)$，得到图式（第一章，3.2.2.7）：
\begin{center}
\resizebox{0.84\textwidth}{!}{%
$\begin{tikzcd}[ampersand replacement=\&,column sep=3.2em,row sep=2.6em]
X \arrow[r,"\dot u"]
  \& Y \arrow[r,"\dot{p}(u)"] \arrow[dr,"\dot v"']
  \& c(\dot u) \arrow[r,"\dot{q}(u)"] \arrow[d,"\dot\mu(S)"]
  \& X[1] \\
\& \& Z \&
\end{tikzcd}$%
}
\end{center}
于是置：

\VerdierSemanticAnchor{ch3d-1-3-2-connecting}
\begin{equation}
\delta(S)=[q(u)][\mu(S)]^{-1}.
\tag{1.3.2.4}
\end{equation}
由 $K(\mathcal A)$ 的三角结构（第一章，3.3.1）和 $D(\mathcal A)$ 的三角结构
（第二章，2.2.6）的定义可知，三角
$(X,Y,Z,[u],[v],\delta(S))$ 是 $D(\mathcal A)$ 中的特异三角。此外，立即
验证态射 $\delta(S)$ 关于 $S$ 是函子性的（(1.3.1)，性质 a)）。因此，我们
赋予典范函子
$\operatorname{comp}(\mathcal A)\longrightarrow D(\mathcal A)$ 一个
$\delta$-函子结构。所得的 $\delta$-函子称为\emph{典范 $\delta$-函子}。
态射 $\delta(S)$ 称为与正合列 $S$ 相联系的\emph{连接态射}。

\VerdierSemanticAnchor{ch3d-1-3-3}
\medskip\noindent\textbf{1.3.3.}\quad
设 $\mathcal A^o$ 为 $\mathcal A$ 的反范畴。对
$\operatorname{comp}(\mathcal A^o)$ 中的每条正合列 $S$，以
$\delta^o(S)$ 表示与之相联系的连接态射。设 $C$ 为一个范畴，$X$ 为 $C$
的一个对象，$u$ 为 $C$ 的一个态射。在下一个命题中，我们滥用记号，以
$X^o$ 和 $u^o$ 表示 $C$ 的反范畴中对应的对象和态射（回忆，$C$ 的反范畴
与 $C$ 具有相同的对象集和态射集）。

\VerdierSemanticAnchor{ch3d-1-3-4}
\medskip\noindent\textbf{命题 1.3.4.}\quad
\textit{a) 对 $\operatorname{comp}(\mathcal A)$ 中的每条正合列 $S$，有：}
\[
\delta(S[1])=-\delta(S)[1].
\]
\noindent\textit{b) 对 $\operatorname{comp}(\mathcal A)$ 中的每条正合列
$S$，考虑到等同 (1.2.8)，有：}
\[
\delta^o(S^o)=\delta(S)^o[1].
\]
\endgroup
